% appendix/questions/removelocking.tex
% mainfile: ../../perfbook.tex
% SPDX-License-Identifier: CC-BY-SA-3.0

\section{为什么不去掉锁？}
\label{sec:app:questions:Why Not Remove Locking?}
%
\epigraph{许多人以为自己在思考，其实只是在重新排列他们的偏见。}
	 {William James}

毫无疑问，许多人把锁视为并行编程的万恶之源，这也不是完全没有道理的。
确实存在一些重要的例子，其中无锁代码比其有锁的对应版本表现好得多，
其中一些在\cref{sec:advsync:Non-Blocking Synchronization}中有所讨论。

然而，无锁算法并不能保证表现良好和具有良好的可扩展性，
正如\cref{fig:count:Atomic Increment Scalability on x86}
（\cpageref{fig:count:Atomic Increment Scalability on x86}）所示。
此外，作为一般规律，算法越复杂，
将锁与选定的无锁技术相结合的优势就越大，即使有大量的硬件支持也是如此，
正如
\IfEbookSize{
\cref{tab:future:Comparison of Locking and HTM,%
tab:future:Comparison of Locking (Augmented by RCU or Hazard Pointers) and HTM}
（\cpageref{tab:future:Comparison of Locking and HTM,%
tab:future:Comparison of Locking (Augmented by RCU or Hazard Pointers) and HTM}）所示。
}{
\cref{tab:future:Comparison of Locking (Plain and Augmented) and HTM}
（\cpageref{tab:future:Comparison of Locking (Plain and Augmented) and HTM}）所示。
}
\Cref{sec:advsync:Non-Blocking Synchronization}
更深入地研究了\IXacrl{nbs}，这是一种流行的无锁方法论。

作为更一般的规律，用简短口号式的方法来进行并行编程不太可能有好结果。
有些人会说这在许多其他领域也是如此，
尤其是那些最显眼地以简短口号为特色的领域。
